\documentclass[UTF8,11pt]{ctexart}
\usepackage[a4paper,margin=24mm]{geometry}
\usepackage{amsmath,amssymb,amsthm,mathtools,mathrsfs}
\usepackage[all]{xy}
\usepackage{xurl,hyperref,enumitem}
\usepackage{tikz}
\usetikzlibrary{arrows.meta,cd,decorations.markings,calc,patterns}
\hypersetup{hidelinks,breaklinks=true}
\setlength{\emergencystretch}{3em}
\allowdisplaybreaks
\newtheorem*{theorem}{Theorem}
\newtheorem*{thm}{Theorem}
\newtheorem*{lemma}{Lemma}
\newtheorem*{proposition}{Proposition}
\newtheorem*{prop}{Proposition}
\newtheorem*{corollary}{Corollary}
\newtheorem*{cor}{Corollary}
\newtheorem*{claim}{Claim}
\theoremstyle{definition}
\newtheorem*{definition}{Definition}
\newtheorem*{defn}{Definition}
\newtheorem*{remark}{Remark}
\newtheorem*{example}{Example}
\newcommand{\R}{\mathbb R}
\newcommand{\C}{\mathbb C}
\newcommand{\Z}{\mathbb Z}
\newcommand{\Q}{\mathbb Q}
\newcommand{\N}{\mathbb N}
\newcommand{\F}{\mathbb F}
\newcommand{\D}{\mathbb D}
\newcommand{\bB}{\mathbb B}
\newcommand{\bC}{\mathbb C}
\newcommand{\bR}{\mathbb R}
\newcommand{\bZ}{\mathbb Z}
\newcommand{\bN}{\mathbb N}
\newcommand{\bQ}{\mathbb Q}
\newcommand{\bD}{\mathbb D}
\newcommand{\bF}{\mathbb F}
\newcommand{\CP}{\mathbb{CP}}
\newcommand{\RP}{\mathbb{RP}}
\newcommand{\sO}{\mathscr O}
\newcommand{\sI}{\mathscr I}
\newcommand{\sF}{\mathscr F}
\newcommand{\sG}{\mathscr G}
\newcommand{\sH}{\mathscr H}
\newcommand{\sR}{\mathscr R}
\newcommand{\sM}{\mathscr M}
\newcommand{\sV}{\mathscr V}
\newcommand{\sU}{\mathscr U}
\DeclareMathOperator{\Hom}{Hom}
\DeclareMathOperator{\Ext}{Ext}
\DeclareMathOperator{\Tor}{Tor}
\DeclareMathOperator{\Spec}{Spec}
\DeclareMathOperator{\Proj}{Proj}
\DeclareMathOperator{\supp}{supp}
\DeclareMathOperator{\im}{im}
\DeclareMathOperator{\id}{id}
\DeclareMathOperator{\Tr}{Tr}
\DeclareMathOperator{\tr}{tr}
\DeclareMathOperator{\rank}{rank}
\DeclareMathOperator{\ord}{ord}
\DeclareMathOperator{\dist}{dist}
\DeclareMathOperator{\End}{End}
\DeclareMathOperator{\Aut}{Aut}
\DeclareMathOperator{\Gal}{Gal}
\DeclareMathOperator{\vol}{vol}
\DeclareMathOperator{\Cl}{Cl}
\DeclareMathOperator{\coker}{coker}
\DeclareMathOperator{\codim}{codim}
\DeclareMathOperator{\divergence}{div}
\DeclareMathOperator{\Ric}{Ric}
\DeclareMathOperator{\grad}{grad}
\DeclareMathOperator{\Hess}{Hess}
\newcommand{\abs}[1]{\left\lvert#1\right\rvert}
\newcommand{\sabs}[1]{\lvert#1\rvert}
\newcommand{\babs}[1]{\bigl\lvert#1\bigr\rvert}
\newcommand{\Babs}[1]{\Bigl\lvert#1\Bigr\rvert}
\newcommand{\bbabs}[1]{\biggl\lvert#1\biggr\rvert}
\newcommand{\BBabs}[1]{\Biggl\lvert#1\Biggr\rvert}
\newcommand{\norm}[1]{\left\lVert#1\right\rVert}
\newcommand{\snorm}[1]{\lVert#1\rVert}
\newcommand{\bnorm}[1]{\bigl\lVert#1\bigr\rVert}
\newcommand{\Bnorm}[1]{\Bigl\lVert#1\Bigr\rVert}
\newcommand{\bbnorm}[1]{\biggl\lVert#1\biggr\rVert}
\newcommand{\BBnorm}[1]{\Biggl\lVert#1\Biggr\rVert}
\newcommand{\widebar}[1]{\overline{#1}}
\newcommand{\nicefrac}[2]{\frac{#1}{#2}}
\newcommand{\linnprod}[2]{\langle#1,#2\rangle}
\newcommand{\myindex}[1]{#1}
\newcommand{\glsadd}[1]{}
\newcommand{\avoidbreak}{}
\newcommand{\sourceref}[1]{\textnormal{[原文：\nolinkurl{#1}]}}
\newcommand{\thmref}[1]{\sourceref{#1}}
\newcommand{\lemmaref}[1]{\sourceref{#1}}
\newcommand{\propref}[1]{\sourceref{#1}}
\newcommand{\corref}[1]{\sourceref{#1}}
\newcommand{\defnref}[1]{\sourceref{#1}}
\newcommand{\exampleref}[1]{\sourceref{#1}}
\newcommand{\exerciseref}[1]{\sourceref{#1}}
\newcommand{\figureref}[1]{\sourceref{#1}}
\newcommand{\sectionref}[1]{\sourceref{#1}}
\newcommand{\appendixref}[1]{\sourceref{#1}}
\newcommand{\stacksref}[2]{\href{https://stacks.math.columbia.edu/tag/#1}{#2 [#1]}}


\newcommand{\E}{\mathbb E}
\newcommand{\Prob}{\mathbb P}
\newcommand{\Reff}{\mathcal R_{\mathrm{eff}}}
\newcommand{\eps}{\varepsilon}
\DeclareMathOperator{\diam}{diam}
\DeclareMathOperator{\Var}{Var}
\DeclareMathOperator{\Crit}{Crit}
\newcommand{\dd}{\,\mathrm d}

\begin{document}
\section*{068\quad 有限etale代数的Galois等变函数重构}
\subsection*{来源与计数目标}
J. S. Milne, \emph{Fields and Galois Theory}, version 5.10, September 2022，Chapter 8。主命题 Proposition 8.18，印刷pp.110--111；支撑8.6--8.7、8.13、8.15--8.17，pp.106--110。访问日期2026-09-28。
\par\url{https://www.jmilne.org/math/CourseNotes/FT.pdf}
\par\textbf{本文件只计一个问题：}构造有限连续 $G=\operatorname{Gal}(F^{\mathrm{sep}}/F)$ 集合上的等变函数代数，证明其为有限etale代数，并重构每个有限etale $F$-代数。
\subsection*{作者命题与人类证明}
\paragraph{Conventions (Chapter 8).}
\noindent\textit{以下背景与记号据所列原文章节摘编；不是逐字引文，也不是新增证明。}\par
Throughout, $F$ is a field, all rings and $F$-algebras are commutative, and unadorned tensor products are over $F$. An $F$-algebra $A$ is finite if it is finite-dimensional as an $F$-vector space; its degree is $[A:F]$.
\begin{definition}[8.5]
An $F$-algebra $A$ is diagonalizable if it is isomorphic to $F^n$ for some $n$, and it is \emph{\'etale} if $L\otimes A$ is diagonalizable for some field $L$ containing $F$.
\end{definition}

\begin{proposition}[8.6]
The following conditions on a finite $F$-algebra $A$ are equivalent:
\begin{enumerate}[label=(\alph*)]
\item $A$ is \'etale;
\item $L\otimes A$ is reduced for all fields $L$ containing $F$;
\item $A$ is a product of separable field extensions of $F$.
\end{enumerate}
\end{proposition}
\begin{proof}
(a)$\Rightarrow$(b). Let $L$ be a field containing $F$. By hypothesis, there exists a field $L'$ containing $F$ such that $L'\otimes A$ is diagonalizable. Let $L''$ be a field containing (copies of) both $L$ and $L'$ (e.g., take $L''$ to be a quotient of $L\otimes L'$ by a maximal ideal). Then $L''\otimes A=L''\otimes_{L'}L'\otimes A$ is diagonalizable, and the map $L\otimes A\to L''\otimes A$ defined by the inclusion $L\to L''$ is injective, and so $L\otimes A$ is reduced.

(b)$\Rightarrow$(c). In particular, $A=A\otimes F$ is reduced, and so it is a finite product of fields (see the above discussion). Suppose that one of the factor fields $F'$ of $A$ is not separable over $F$. Then $F$ has characteristic $p\ne0$ and there exists an element $u$ of $F'$ whose minimal polynomial is of the form $g(X^p)$ with $g\in F[X]$ (see 3.6 et seq.). Let $L$ be a field containing $F$. Then
\[L\otimes F[u]\simeq L\otimes(F[X]/(g(X^p)))\simeq L[X]/(g(X^p)).\]
If $L$ is chosen so that the coefficients of $g(X)$ become $p$th powers in it, then $g(X^p)$ is a $p$th power in $L[X]$ (see the proof of 2.24), and so $L\otimes F[u]$ is not reduced. But $L\otimes F[u]\subset L\otimes A$, and so this contradicts the hypothesis.

(c)$\Rightarrow$(a). We may suppose that $A$ itself is a separable field extension of $F$. From the primitive element theorem (5.1), we know that $A=F[u]$ for some $u$. Because $F[u]$ is separable over $F$, the minimal polynomial $f(X)$ of $u$ is separable, which means that, in any splitting field $L$ for $f$,
\[f(X)=\prod(X-u_i),\qquad u_i\ne u_j\text{ for }i\ne j.\]
Now $L\otimes A\simeq L\otimes F[X]/(f)\simeq L[X]/(f)$, and, according to the Chinese remainder theorem (8.1),
\[L[X]/(f)\simeq\prod_i L[X]/(X-u_i)\simeq L\times\cdots\times L.\]
\end{proof}
\begin{corollary}[8.7]
An $F$-algebra $A$ is \'etale if and only if $F^{\mathrm{sep}}\otimes A$ is diagonalizable.
\end{corollary}
\begin{proof}
The proof that (c) implies (a) in (8.6) shows that $L\otimes A$ is diagonalizable if certain separable polynomials split in $L$. By definition, all separable polynomials split in $F^{\mathrm{sep}}$.
\end{proof}

\paragraph{8.13: Galois descent for vector spaces.}
Let $\Omega$ be a Galois extension of $F$ with Galois group $G$, and let $V$ be a finite-dimensional vector space over $F$. Write $V_\Omega=\Omega\otimes_F V$. The group $G$ acts on $V_\Omega$ through its action on $\Omega$. The map
\[v\longmapsto1\otimes v:V\longrightarrow (V_\Omega)^G\]
is an isomorphism. To see this, choose an $F$-basis $\{e_1,\ldots,e_n\}$ for $V$. Then $\{e_1,\ldots,e_n\}$ is also an $\Omega$-basis for $V_\Omega$, and
\[\sigma(a_1e_1+\cdots+a_ne_n)=(\sigma a_1)e_1+\cdots+(\sigma a_n)e_n,\qquad a_i\in\Omega.\]
Therefore $a_1e_1+\cdots+a_ne_n$ is fixed by all $\sigma\in G$ if and only if $a_1,\ldots,a_n\in F$.

\paragraph{The two constructions (8.14--8.16).}
We fix a separable closure $\Omega$ of $F$, and let $G=\operatorname{Gal}(\Omega/F)$. By a $G$-set we mean a set $S$ equipped with an action of $G$ such that $G\times S\to S$ is continuous with respect to the Krull topology on $G$ and the discrete topology on $S$.

For an \'etale $F$-algebra $A$, let $\mathcal F(A)$ denote the set of $F$-algebra homomorphisms $f:A\to\Omega$. We let $G$ act on $\mathcal F(A)$ through its action on $\Omega$,
\[(\sigma f)(a)=\sigma(f(a)).\]
For some finite Galois extension $E$ of $F$ in $\Omega$, the images of all homomorphism $A\to\Omega$ are contained in $E$, and so the action of $G$ on $\mathcal F(A)$ factors through $\operatorname{Gal}(E/F)$. Therefore $\mathcal F(A)$ is a $G$-set.

Let $A=A_1\times\cdots\times A_n$ with each $A_i$ an \'etale $F$-algebra. Because $\Omega$ is an integral domain, every homomorphism $f:A\to\Omega$ is zero on all but one $A_i$, and so, to give a homomorphism $A\to\Omega$ amounts to giving a homomorphism $A_i\to\Omega$ for some $i$. In other words,
\[\mathcal F\left(\prod_i A_i\right)\simeq\coprod_i\mathcal F(A_i).\]
In particular, for an \'etale $F$-algebra $A=\prod_iF_i$, $F_i$ a field,
\[\mathcal F(A)\simeq\coprod_i\operatorname{Hom}_F(F_i,\Omega).\]
From Proposition 2.12, we deduce that $\mathcal F(A)$ is finite of order $[A:F]$.

For a $G$-set $S$, we let $G$ act on the $F$-algebra $\operatorname{Hom}(S,\Omega)$ of maps $S\to\Omega$ through its actions on $S$ and $\Omega$,
\[(\sigma f)(s)=\sigma(f(\sigma^{-1}s)).\]
We define $\mathcal A(S)$ to be the set of elements of $\operatorname{Hom}(S,\Omega)$ fixed by $G$. Thus $\mathcal A(S)$ is the $F$-subalgebra consisting of the maps $f:S\to\Omega$ such that $f(\sigma s)=\sigma f(s)$ for all $\sigma\in G$, $s\in S$.

Suppose that $G$ acts transitively on $S$. Choose an $s\in S$, and let $H\subset G$ be its stabilizer. Then $H$ is an open subgroup of $G$, and so $E=\Omega^H$ is a finite extension of $F$ (7.13). An element $f$ of $\mathcal A(S)$ is determined by its value on $s$, which can be any element of $\Omega$ fixed by $H$. It follows that
\[f\longmapsto f(s):\mathcal A(S)\longrightarrow E\]
is an isomorphism of $F$-algebras.

Every element of $S$ is of the form $\sigma s$ with $\sigma\in G$, and $\sigma s=\sigma's$ if and only if $\sigma H=\sigma'H$. Similarly, every element of $\mathcal F(E)$ is of the form $\sigma|_E$ with $\sigma\in G$, and $\sigma|_E=\sigma'|_E$ if and only if $\sigma H=\sigma'H$. It follows that the map
\[\sigma s\longmapsto\sigma|_E:S\longrightarrow\mathcal F(E)\]
is an isomorphism of $G$-sets.

Let $E$ be a finite separable extension of $F$. Let $S=\mathcal F(E)$ and choose an $s\in S$, i.e., an embedding $s:E\hookrightarrow\Omega$. The above calculation shows that $\mathcal A(S)=sE$. In particular, $s$ defines an isomorphism $E\to\mathcal A(\mathcal F(E))$.

\begin{proposition}[8.17]
Let $S$ be a finite $G$-set, and let $S=S_1\sqcup\cdots\sqcup S_n$ be its decomposition into $G$-orbits. For each $i$, choose an $s_i\in S_i$, and let $F_i$ be the subfield of $\Omega$ fixed by the stabilizer of $s_i$.
\begin{enumerate}[label=(\alph*)]
\item Each $F_i$ is a finite separable extension of $F$.
\item $f\mapsto(f(s_1),\ldots,f(s_n))$ defines an isomorphism $\mathcal A(S)\to F_1\times\cdots\times F_n$ of $F$-algebras.
\item The map sending $\sigma s_i\in S_i\subset S$ to $\sigma|_{F_i}\in\mathcal F(F_i)\subset\mathcal F(F_1\times\cdots\times F_n)$ is an isomorphism of $G$-sets $S\to\mathcal F(F_1\times\cdots\times F_n)$.
\end{enumerate}
\end{proposition}
\begin{proof}
This follows directly from the special case considered in 8.16.
\end{proof}
\begin{proposition}[8.18]
For every finite $G$-set $S$, the $F$-algebra $\mathcal A(S)$ is \'etale with degree equal to $|S|$. Moreover, every \'etale $F$-algebra $A$ is of the form $\mathcal A(S)$ for some $G$-set $S$. More precisely, $A\simeq\mathcal A(\mathcal F(A))$.
\end{proposition}
\begin{proof}
The first statement follows from (8.6) and (8.17). We prove the third statement. There is a canonical isomorphism of $\Omega$-algebras
\[a\otimes c\longmapsto(\sigma a\cdot c)_{\sigma\in\mathcal F(A)}:\Omega\otimes A\longrightarrow\prod_{\sigma\in\mathcal F(A)}\Omega.\]
When we let $G$ act on $\Omega\otimes A$ through $\Omega$, and pass to the fixed elements, we obtain an isomorphism
\[A\overset{8.13}{=}(\Omega\otimes A)^G\simeq\mathcal A(\mathcal F(A)).\]
This implies the second statement of the proposition (which can also be deduced from 8.17).
\end{proof}
\subsection*{整理说明（非作者正文）}
本条锁定作者8.18的双向对象重构结论，未附加8.21的整个范畴反变等价，也没有把8.19留作练习的求值映射唯一性写成已经取得的证明。核心轨道--稳定子--不动域对应、代数的分裂和取不变量均已展开。8.17引用的8.16论证完整收录，不只保留最后三行证明。

标准先修为有限约化代数的中国剩余分解、域扩张的原始元定理、可分嵌入计数、无限Galois理论中开子群与有限子扩张对应。8.18的原文将 $a\otimes c$ 的映射写在 $\Omega\otimes A$ 上；此处按原文保留因子顺序，读者可用张量积的交换同构对应记号，不静默改写成另一原句。

难度依据：需要把有限群作用数据转成代数并检验反向重构，不是只将子群换名为中间域。来自作者PDF可提取正文，数学符号转为TeX；前置记号按原章摘编，具体环境标题为排版加入。未取得PDF本地副本，保存的是转录摘录。
\subsection*{可修改的简短标注（非作者正文）}
$c$：有限etale代数的分类与重构。

$a$：Galois连续作用；轨道与稳定子；不动域；等变函数代数；张量分裂；中国剩余定理；不变量下降。

\texttt{cross\_theory\_status = yes}。以有限连续作用的轨道确定有限可分域，再将等变函数代数与域的直积识别；反向把分裂代数取Galois不变量以恢复原代数。位置：8.16--8.18。
标签为非穷尽、可修改初稿。
\subsection*{许可与修改记录}
原作者及作品见来源栏；来源采用 CC BY-NC-SA 4.0。本题证摘录和附加整理沿用该许可。2026-09-28：增加题号、来源定位、独立排版、整理说明和初稿标注；数学内容的取材范围及排版变化见整理说明。
\end{document}
